\documentclass[12pt]{article}
\usepackage{amsmath}
\usepackage{amssymb}
\usepackage{geometry}
\usepackage{ulem}

\geometry{margin=2.5cm}

\begin{document}

\section{Proving Q}

\noindent To derive the risk-neutral probability $q$, we consider an investor who allocates 
wealth between the stock $S$ and the bond $B$ by choosing a portfolio 
$H = (\alpha, \beta)$. The portfolio value at $t=1$ is:

\vspace{1em}
\[
V_1 = \alpha S_1 + \beta \cdot B_1 =
\left\{
\begin{array}{l}
\alpha \cdot S_0 \cdot u + \beta B_1 = X_{\text{up}} \\[4pt]
\alpha \cdot S_0 \cdot d + \beta B_1 = X_{\text{down}}
\end{array}
\right.
\]
\vspace{1em}

\noindent The investor maximises expected utility $\mathbb{E}^P[u(V_1)]$ subject to the budget 
constraint $\alpha S_0 + \beta B_0 = V_0$. We show that the first-order conditions 
uniquely pin down a probability measure $q$ (the \textbf{risk-neutral probability}) under which all assets earn the risk-free rate $r$.

\vspace{2.5em}

\noindent \uline{The investor's optimisation problem is thus:}

\[
\max_{\alpha,\, \beta} \; \mathbb{E}^P\left[u\!\left(\alpha S_1 + \beta \cdot B_1\right)\right]
\quad \text{s.t.} \quad
\alpha S_0 + \beta B_0 = V_0
\]
\vspace{0.1em}

\noindent {\small where the constraint states that the cost of the portfolio today equals
initial wealth $V_0$. Since the constraint is an equality, we maximise
using the Lagrangian:}
\vspace{0.1em}

\begin{align*}
\mathcal{L}(\alpha, \beta, \lambda)
  &= \mathbb{E}^P\left[u\!\left(\alpha S_1 + \beta B_1\right)\right]
     - \lambda\!\left(\alpha S_0 + \beta B_0 - V_0\right) \\[6pt]
  &= p \cdot u(X_{\text{up}}) + (1-p) \cdot u(X_{\text{down}})
     - \lambda\!\left(\alpha S_0 + \beta B_0 - V_0\right) \\[6pt]
  &= p \cdot u(\alpha S_0 \cdot u + \beta B_1)
     + (1-p) \cdot u(\alpha S_0 \cdot d + \beta B_1)
     - \lambda\!\left(\alpha S_0 + \beta B_0 - V_0\right)
\end{align*}

\vspace{1em}

\noindent Taking first-order conditions with respect to $\alpha$ and $\beta$:
\begin{align*}
\mathcal{L}_\alpha
  &= p \cdot u'(X_{\text{up}}) \cdot S_0 \cdot u + (1-p) \cdot u'(X_{\text{down}}) \cdot S_0 \cdot d
     - \lambda \cdot S_0 = 0
  \\[4pt]
  &= \mathbb{E}^P\left[u'(V_1) \cdot S_1\right] = \lambda S_0
\end{align*}

\noindent {\small Isolating $S_0$:}
\[
\longrightarrow \quad
S_0 = \frac{1}{\lambda} \cdot \mathbb{E}^P\left[u'(V_1) \cdot S_1\right]
\tag{I}
\]

\vspace{2em}

\begin{align*}
\mathcal{L}_\beta
  &= p \cdot u'(X_{\text{up}}) \cdot B_1 + (1-p) \cdot u'(X_{\text{down}}) \cdot B_1
     - \lambda \cdot B_0 = 0
  \\[4pt]
  &= \mathbb{E}^P\left[u'(V_1) \cdot B_1\right] = \lambda B_0
\end{align*} 

\noindent {\small Isolating $\lambda$:}
\[
\longrightarrow \quad
\lambda = \frac{\mathbb{E}^P\left[u'(V_1) \cdot B_1\right]}{B_0}
\tag{II}
\]

\noindent We normalize the bond price to $B_0 = 1$ (one unit invested today) 
and $B_1 = (1+r)$ (grows at the risk-free rate):

\[
\frac{1}{\lambda}
  = \frac{1}{(1+r) \cdot \mathbb{E}^P\left[u'(V_1)\right]}
\]
\vspace{0.1em}

\noindent {\small Substituting this expression for $\frac{1}{\lambda}$ into (I):}
\begin{align*}
S_0
  &= \frac{1}{1+r} \cdot \frac{\mathbb{E}^P\left[u'(V_1) \cdot S_1\right]}{\mathbb{E}^P\left[u'(V_1)\right]}
  \tag{III} \\[4pt]
\end{align*}

\noindent \textbf{\underline{We define the variable $Z$}}, which represents the marginal utility of $V_1$ divided by the
expected value of the marginal utility of $V_1$:
\[
Z = \frac{u'(V_1)}{E^P\!\left[u'(V_1)\right]}
\]
\vspace{0.5em}

\noindent Since $V_1$ takes two values in the binomial model, 
$Z$ takes two values accordingly:
\[
Z_{\text{up}} =
  \frac{u'(X_{\text{up}})}
       {p \cdot u'(X_{\text{up}}) + (1-p) \cdot u'(X_{\text{down}})}
\qquad\qquad
Z_{\text{down}} =
  \frac{u'(X_{\text{down}})}
       {p \cdot u'(X_{\text{up}}) + (1-p) \cdot u'(X_{\text{down}})}
\]

\vspace{1.5em}

\noindent\uline{$Z$ satisfies two properties:}

\begin{enumerate}
  \item Since utility is increasing, $u'(V_1) > 0$, so $Z > 0$ 
        in both states.

  \vspace{0.2em}

  \item The expected value of $Z$ is equal to 1 by construction:
        \[
          E^P[Z]
          = E^P\!\left[\frac{u'(V_1)}{E^P\!\left[u'(V_1)\right]}\right]
          = \frac{E^P\!\left[u'(V_1)\right]}{E^P\!\left[u'(V_1)\right]}
          = 1
        \]
\end{enumerate}

\vspace{1.5em}

\noindent {\small Using the two properties of $Z$, we can rewrite equation (III) as:}
\[
S_0 = \frac{1}{1+r} \cdot \mathbb{E}^P\!\left[Z \cdot S_1\right]
\]

\noindent {\small Expanding $\mathbb{E}^P[Z \cdot S_1]$ explicitly in the 
binomial model, where $S_1 = S_0 u$ in the up state and $S_1 = S_0 d$ 
in the down state:}
\begin{align*}
 & \mathbb{E}^P\!\left[Z \cdot S_1\right] 
  = p \cdot Z_{\text{up}} \cdot S_0 \cdot u + (1-p) \cdot Z_{\text{down}} \cdot S_0 \cdot d \\[0.5pt]
\\[-8pt]
\longrightarrow \quad & S_0 = \frac{1}{1+r} \cdot \left( p \cdot Z_{\text{up}} \cdot S_0 \cdot u + (1-p) \cdot Z_{\text{down}} \cdot S_0 \cdot d \right)
\end{align*}
\vspace{1em}

\noindent {\small By dividing both sides by $S_0$, we get:}
\[
1 + r = p \cdot Z_{\text{up}} \cdot u + (1-p) \cdot Z_{\text{down}} \cdot d
\tag{IV}
\]

\noindent We define $q$ and $(1-q)$ as the real-world probability reweighted by how much the investor values a dollar in that state, which makes the investor indifferent between the stock and the risk-free bond:
\[
q = p \cdot Z_{\text{up}}
  = p \cdot \frac{u'(X_{\text{up}})}{E\!\left[u'(V_1)\right]}
\qquad \text{and} \qquad
1-q = (1-p) \cdot Z_{\text{down}}
       = (1-p) \cdot \frac{u'(X_{\text{down}})}{E\!\left[u'(V_1)\right]}
\]

\vspace{1em}

\noindent This is valid probability measure because both terms 
are strictly positive (since $Z > 0$, $p \in (0,1)$), and they sum to 1:
\[
q + (1-q) 
  = p \cdot Z_{\text{up}} + (1-p) \cdot Z_{\text{down}} 
  = \mathbb{E}^P[Z] = 1 \quad 
\]
\vspace{0.1em}

\noindent {\small Substituting into equation (IV) and solving for $q$:}
\begin{align*}
1 + r &= q \cdot u + (1-q) \cdot d \\[4pt]
1 + r &= q(u - d) + d
\end{align*}

\[
\boxed{q = \frac{(1+r) - d}{u - d}}
\]
\vspace{1em}

\noindent This is the \textbf{risk-neutral probability}: the unique 
value of $q$ consistent with no-arbitrage, independent of any investor's 
preferences. 
\end{document}